% ============================================================
% INTRODUCTION (4 paragraphs: problem, what is known, the gap, the solution + H1/H2)
% ============================================================
\section{Introduction}
\label{sec:intro}

% --- Paragraph 1: the basic problem ---
Pain has no objective marker, so its assessment depends on how a clinician
interprets a patient's report and behaviour, and that interpretation is not
neutral. Racial and ethnic minority patients are undertreated for pain relative
to non-Hispanic White patients across experimental, post-operative, cancer and
chronic pain and in the emergency department \cite{green2003unequal,
anderson2009racial}, a disparity that a meta-analysis of twenty years of United
States evidence found across settings and pain types \cite{meghani2012time} and
that emergency-department data confirmed for opioid prescribing
\cite{pletcher2008trends}. The disparity has a cognitive substrate: half of a
sample of White medical students and residents endorsed false beliefs about
biological differences between Black and White people, and those who did rated a
Black patient's pain lower \cite{hoffman2016racial}. Large language models
(LLMs) now encode clinical knowledge at a level that invites deployment
\cite{singhal2023large, thirunavukarasu2023large}, are being evaluated across a
widening range of clinical tasks \cite{bedi2025testing}, and are being embedded
in documentation, triage and decision-support tools. A commercial risk algorithm
has already been shown to reproduce racial bias at the scale of millions of
patients \cite{obermeyer2019dissecting}. A language model that learned its
clinical judgement from text produced by a system with documented disparities
may carry those disparities into every encounter it touches, at scale and
without the variability that lets an individual clinician's bias be noticed.

% --- Paragraph 2: what is known ---
The evidence on whether LLMs do so is mixed and tracks model generation. Q-Pain,
a question-answering dataset for social bias in pain management, found
statistically significant differences in treatment between intersectional
race--gender subgroups for GPT-2 and GPT-3 \cite{loge2021qpain}. Commercial
models of 2023 reproduced race-based medical misconceptions and answered the same
question inconsistently on repetition \cite{omiye2023race}; GPT-4 generated
demographically stereotyped clinical vignettes \cite{zack2024gpt4}; GPT-3.5 and
GPT-4 projected higher costs and longer hospitalisations for White patients from
real cases \cite{yang2024unmasking}; and systematic and scoping reviews of 38 and
92 studies respectively describe bias as pervasive across clinical applications
\cite{suenghataiphorn2025systematic, xiao2026scoping}. Pain is the exception.
When Young et al.\ presented GPT-4 and Gemini with 40 real cases from MIMIC-IV
in every race/ethnicity and sex combination, race/ethnicity and sex did not
influence pain ratings or opioid recommendations \cite{young2025racial}.
Counterfactual audits of emergency triage, meanwhile, have found large
gender-swap flip rates in proprietary models \cite{young2026equitriage},
sensitivity to stigmatising language in a single sentence
\cite{huang2026artificial}, and, in an outcome-grounded design, harmful
reprioritisation from psychiatric history, homelessness and substance use but
not from race, language or insurance status \cite{jarrett2026outcome}. Two
methodological results frame all of these: a model changes its own answer when
the prompt is re-run unchanged, even under nominally deterministic settings
\cite{atil2024nondeterminism, song2025good}, and that instability floor must be
measured before a counterfactual difference can be read as bias
\cite{bellibatlu2026instability}. The multilingual PAINT resource states the
question that remains open: whether LLMs reproduce known disparities in pain
assessment \cite{zhang2026paint}.

% --- Paragraph 3: the gap (the hook) ---
Four gaps limit what that literature can say. First, every pain audit tested
proprietary models reached through an API, whereas hospitals increasingly run
open-weight models on their own hardware for privacy and cost
\cite{nowak2026secure}; open-source models have been audited for triage
\cite{aljohani2026counterfactual} and diagnosis \cite{ahmed2026auditing}, but
not for pain. Second, every pain audit announced the patient's race and sex
explicitly. Instruction-tuned models are trained to respond carefully to
protected attributes \cite{ouyang2022training}, so an explicit label may trigger
a guarded response that an indirect cue would not, and different cues for the
same group have been shown to yield inconsistent conclusions about bias, in both
magnitude and direction \cite{tonneau2026cues}. In real notes, demographic
signal often arrives through a name, and name-based audits outside medicine
consistently find that models disadvantage names associated with racial
minorities and women \cite{salinas2024name, an2024hiring, nghiem2024gotta,
bhetwal2026name}, as human employers did in the original correspondence
audits \cite{bertrand2004emily}. No pain audit has used name cues. Third, where
names have been used, one name per group leaves a name's idiosyncrasy
indistinguishable from a group effect, and a recent preregistered audit found
that the audit instrument itself can produce effects as large as the demographic
bias under test \cite{vohra2026audit}. Fourth, few audits report how often a
model changes its own answer when nothing changes, so a measured difference has
no baseline against which to be judged.

% --- Paragraph 4: how this study addresses it; purposes and hypotheses ---
This study addresses the first three gaps and supplies the fourth with
counterfactual vignette audits \cite{kusner2017counterfactual} that vary four
races and two sexes while holding every other token of the prompt fixed: of
four open-weight models on six acute-pain cases (Study~1), and of a wider panel
of local and API-served models on those cases and on the 50 clinician-written
Q-Pain cases (Study~2). The primary purpose was to
determine whether open-weight models reproduce the human disparity when race and
sex are stated explicitly. Hypothesis 1 (H1), following the null result of
Young et al.\ \cite{young2025racial}, predicted that explicit labels would not
change pain-severity ratings or strong-opioid recommendations. The secondary
purpose was to test whether indirect cues reveal a disparity that explicit
labels mask. Hypothesis 2 (H2) predicted that when identity was cued only by the
patient's name, models would rate pain lower and recommend strong opioids less
often for patients with non-White names, and with Black and Hispanic names in
particular, than for patients with White male names. No hypothesis concerned sex,
because the human evidence on sex is less uniform than on race
(Section~\ref{sec:related}); the woman minus man contrasts are reported as
exploratory, with the human finding of less opioid analgesia for women in emergency
care as the direction to watch. A third purpose, addressed
by Study~2 without a directional hypothesis, was to establish whether the answer
depends on the cases or the models chosen: whether it holds on clinician-written
cases, for proprietary models and for open-weight models from 3 billion parameters to
mixture-of-experts models of several hundred billion, with reasoning switched on as well as off, and at 16-bit as well as
4-bit precision. Every comparison was paired on the same case and sampling seed,
tested with bootstrap intervals and permutation tests, corrected for multiple
comparisons, and set against the model's own re-sampling noise floor.
